\begin{frame}{그림으로: 데이터 분포 $\leftrightarrow$ 노이즈}
  \begin{center}
  \includegraphics[height=5.4cm]{gen_diffusion_forward_reverse.png}
  \end{center}
  \vskip0.2em
  \small
  정방향(위쪽 화살표, 고정)은 두 봉우리 데이터 분포를 표준정규
  $\mathcal{N}(0,1)$까지 점점 뭉갠다. 역방향(아래쪽 화살표,
  $\epsilon_\theta$가 학습)은 그 과정을 거슬러 \textbf{노이즈에서
  데이터 분포를 되살린다}.
\end{frame}
